\subsubsection{Interpretable Deep Convolutional Model for Nonlinear Multivariate Time Series in Complex Systems}

\citet{baric2026dcits} propose the Deep Convolutional Interpreter for Time
Series (DCIts), an intrinsically interpretable model for one-step multivariate
forecasting. The paper continues the benchmark of
\citet{baric2021benchmarking}: it reuses Datasets 1--8, modifies the switching
process in Dataset 8, and adds controlled VAR(2) and cubic processes. In all
these experiments, the equations used to generate the observations supply the
reference target--source--lag structure and, where applicable, the exact signed
coefficients, biases, regimes, and polynomial orders.

\begin{landscape}
\begingroup
\scriptsize
\setlength{\LTpre}{0.4\baselineskip}
\setlength{\LTpost}{0.4\baselineskip}
\setlength{\tabcolsep}{4pt}
\renewcommand{\arraystretch}{1.08}
\begin{longtable}{@{}
>{\raggedright\arraybackslash}p{0.10\linewidth}
>{\raggedright\arraybackslash}p{0.49\linewidth}
>{\raggedright\arraybackslash}p{0.35\linewidth}@{}}
\caption{Exact transcription of Table~III in
\citet{baric2026dcits}: benchmarking datasets used in the study.}
\label{tab:dcits-datasets-exact}\\
\toprule
Name & Model & Parameters \\
\midrule
\endfirsthead
\multicolumn{3}{@{}l}{\small\emph{Table~\thetable\ continued}}\\
\toprule
Name & Model & Parameters \\
\midrule
\endhead
\midrule
\multicolumn{3}{r@{}}{\small\emph{Continued on the next page}}\\
\endfoot
\bottomrule
\endlastfoot

Dataset 1 &
$X_{n,t}=\alpha_n^{\mathrm{gt}}+\varepsilon_{n,t}$ &
$\alpha_n^{\mathrm{gt}}\sim\mathcal{N}(0,1)$ \\

Dataset 2 &
$X_{n,t}=\alpha_{n,n,l}^{\mathrm{gt}}X_{n,t-l}+\varepsilon_{n,t}$ &
$\alpha_{n,n,3}^{\mathrm{gt}}=\alpha_{n,n,7}^{\mathrm{gt}}=1/2$ \\

Dataset 3 &
$X_{n,t}=\tanh\!\left(\alpha_{n,n,l}^{\mathrm{gt}}X_{n,t-l}\right)
+\varepsilon_{n,t}$ &
$\alpha_{n,n,3}^{\mathrm{gt}}=5/7,\quad
 \alpha_{n,n,7}^{\mathrm{gt}}=\alpha_{n,n,9}^{\mathrm{gt}}=1/7$ \\

Dataset 4 &
$\begin{aligned}
X_{1,t}&=\alpha_{1,2,l}^{\mathrm{gt}}X_{2,t-l}+\varepsilon_{1,t},\\
X_{2,t}&=\alpha_{2,1,l}^{\mathrm{gt}}X_{1,t-l}+\varepsilon_{2,t}
\end{aligned}$ &
$\begin{aligned}
\alpha_{1,2,2}^{\mathrm{gt}}&=\alpha_{1,2,9}^{\mathrm{gt}}=2/5,
&\alpha_{1,2,5}^{\mathrm{gt}}&=1/5,\\
\alpha_{2,1,2}^{\mathrm{gt}}&=\alpha_{2,1,9}^{\mathrm{gt}}=2/5,
&\alpha_{2,1,5}^{\mathrm{gt}}&=1/5
\end{aligned}$ \\

Dataset 5 &
$X_{n,t}=\alpha_{n,1,l}^{\mathrm{gt}}X_{1,t-l}+\varepsilon_{n,t}$ &
$\begin{aligned}
\alpha_{1,1,3}^{\mathrm{gt}}&=1/2,&
\alpha_{1,1,4}^{\mathrm{gt}}&=1/2,&
\alpha_{2,1,9}^{\mathrm{gt}}&=1,\\
\alpha_{3,1,2}^{\mathrm{gt}}&=1/2,&
\alpha_{3,1,7}^{\mathrm{gt}}&=1/2,\\
\alpha_{4,1,3}^{\mathrm{gt}}&=1/10,&
\alpha_{4,1,4}^{\mathrm{gt}}&=1/10,&
\alpha_{4,1,8}^{\mathrm{gt}}&=4/5,\\
\alpha_{5,1,2}^{\mathrm{gt}}&=1/3,&
\alpha_{5,1,5}^{\mathrm{gt}}&=2/9,&
\alpha_{5,1,8}^{\mathrm{gt}}&=4/9
\end{aligned}$ \\

Dataset 6 &
$X_{n,t}=\tanh\!\left(\alpha_{n,1,l}^{\mathrm{gt}}X_{1,t-l}\right)
+\varepsilon_{n,t}$ &
Identical to Dataset 5 \\

Dataset 7 &
$\begin{aligned}
X_{1,t}&=\alpha_{1,1,1}^{\mathrm{gt}}X_{1,t-1}
       +\alpha_{1,1,5}^{\mathrm{gt}}X_{1,t-5}+\varepsilon_{1,t},\\
X_{2,t}&=1+\alpha_{2,1,2}^{\mathrm{gt}}X_{1,t-2}+\varepsilon_{2,t},\\
X_{3,t}&=\alpha_{3,2,1}^{\mathrm{gt}}X_{2,t-1}
       +\alpha_{3,4,4}^{\mathrm{gt}}X_{4,t-4}+\varepsilon_{3,t},\\
X_{4,t}&=1+\alpha_{4,3,4}^{\mathrm{gt}}X_{3,t-4}
       +\alpha_{4,5,1}^{\mathrm{gt}}X_{5,t-1}+\varepsilon_{4,t},\\
X_{5,t}&=\alpha_{5,5,4}^{\mathrm{gt}}X_{5,t-4}
       +\alpha_{5,2,1}^{\mathrm{gt}}X_{2,t-1}+\varepsilon_{5,t}
\end{aligned}$ &
$\begin{aligned}
\alpha_{1,1,1}^{\mathrm{gt}}&=1/4,&
\alpha_{1,1,5}^{\mathrm{gt}}&=3/4,\\
\alpha_{2,1,2}^{\mathrm{gt}}&=-1,\\
\alpha_{3,2,1}^{\mathrm{gt}}&=1,&
\alpha_{3,4,4}^{\mathrm{gt}}&=1,\\
\alpha_{4,3,4}^{\mathrm{gt}}&=-2/7,&
\alpha_{4,5,1}^{\mathrm{gt}}&=-5/7,\\
\alpha_{5,5,4}^{\mathrm{gt}}&=12/22,&
\alpha_{5,2,1}^{\mathrm{gt}}&=10/22
\end{aligned}$ \\

Dataset 8 &
$\begin{aligned}
&\text{if }X_{1,t-5}>1/2:\\
&X_{1,t}=b_{\mathrm{high}}+\varepsilon_{1,t},\\
&X_{2,t}=\alpha_{2,1,5}^{\mathrm{gt}}X_{1,t-5}+\varepsilon_{2,t},\\
&X_{3,t}=\alpha_{3,1,4}^{\mathrm{gt}}X_{1,t-4}+\varepsilon_{3,t},\\
&X_{4,t}=\alpha_{4,4,1}^{\mathrm{gt}}X_{4,t-1}
       +\alpha_{4,4,4}^{\mathrm{gt}}X_{4,t-4}+\varepsilon_{4,t},\\
&\text{else:}\\
&X_{1,t}=b_{\mathrm{low}}+\varepsilon_{1,t},\\
&X_{2,t}=\alpha_{2,4,2}^{\mathrm{gt}}X_{4,t-2}+\varepsilon_{2,t},\\
&X_{3,t}=\alpha_{3,4,4}^{\mathrm{gt}}X_{4,t-4}+\varepsilon_{3,t},\\
&X_{4,t}=\alpha_{4,4,1}^{\mathrm{gt}}X_{4,t-1}
       +\alpha_{4,4,4}^{\mathrm{gt}}X_{4,t-4}+\varepsilon_{4,t}
\end{aligned}$ &
$\begin{aligned}
\alpha_{2,1,5}^{\mathrm{gt}}&=4/5,&
\alpha_{2,4,2}^{\mathrm{gt}}&=2/3,\\
\alpha_{3,1,4}^{\mathrm{gt}}&=2/3,&
\alpha_{3,4,4}^{\mathrm{gt}}&=4/5,\\
\alpha_{4,4,1}^{\mathrm{gt}}&=1/2,&
\alpha_{4,4,4}^{\mathrm{gt}}&=2/5,\\
b_{\mathrm{high}}&=7/10,&
b_{\mathrm{low}}&=2/10
\end{aligned}$

\medskip
$b_{\mathrm{high}}$ and $b_{\mathrm{low}}$ switch, based on a persistence
duration, $P$.

\medskip
$P\sim\max\!\left(P_{\min},\mathcal{N}(\mu_P,\sigma_P)\right)$ \\

\end{longtable}
\endgroup
\end{landscape}

\paragraph{Reading the original benchmarking-dataset table.}
Table~\ref{tab:dcits-datasets-exact} above transcribes Table~III in Appendix~A
of \citet{baric2026dcits}. The DCIts paper presents these as the corrected
specifications of Datasets~1--8 and explicitly modifies the switching process
of Dataset~8 relative to \citet{baric2021benchmarking}. The notation
$\alpha^{\mathrm{gt}}$ is retained because the table uses each generator
coefficient directly as an explanation ground truth.

For an input window $Q_t\in\mathbb{R}^{N\times L}$, DCIts predicts target series
$k$ one step ahead through the following re-indexed form of the paper's
prediction equation:

\begin{equation}
    \widehat{x}_{k,t+1}
    = \sum_{j=1}^{N}\sum_{\ell=1}^{L}
    (\alpha^{\mathrm{DCI}}_t)_{k,j,\ell}\,(Q_t)_{j,\ell},
\end{equation}

where $(\alpha^{\mathrm{DCI}}_t)_{k,j,\ell}$ is a sample-specific coefficient
linking source series $j$ at lag $\ell$ to target series $k$. DCIts orders the
window as $Q_t=(X_t,X_{t-1},\ldots,X_{t-L+1})$. Accordingly,
$(Q_t)_{j,1}=x_{j,t}$ is the latest observed input and the forecast is
$x_{k,t+1}$. On the common axis this is $\tau=t$: the latest input is relative
time 0 and the one-step target is relative time $+1$. The model factorizes
this transition tensor as

\begin{equation}
    \alpha^{\mathrm{DCI}}_t = C_t \odot F_t,
\end{equation}

where the Focuser $F_t\in(0,1)^{N\times N\times L}$ softly selects relevant
source--lag entries and the Modeler $C_t\in\mathbb{R}^{N\times N\times L}$
assigns their signed local magnitudes. The coefficient
$(\alpha^{\mathrm{DCI}}_t)_{k,j,\ell}$ describes the learned local interaction,
while $(\alpha^{\mathrm{DCI}}_t)_{k,j,\ell}\,(Q_t)_{j,\ell}$ is the realized
numerical contribution of that entry to the current forecast. Thus,
coefficient recovery evaluates the learned transition mechanism, and the
coefficient--input product explains how that mechanism contributes to one
particular output.

The generator equations permit evaluation of support, lag, sign, and magnitude.
For example, if the generator contains
$x_{2,t}=1-x_{1,t-2}+\varepsilon_{2,t}$, the reference specifies target
$x_2$, source $x_1$, lag 2, coefficient $-1$, and bias $1$.

\paragraph{Evaluation calculations and reported comparisons.}
DCIts first checks forecast quality with test-set MSE. Table~II of
\citet{baric2026dcits} reports the mean and standard deviation of MSE over
repeated training runs for DCIts and IMV-LSTM. Figure~3 displays the
coefficient of variation
$\operatorname{sd}(\mathrm{MSE})/\operatorname{mean}(\mathrm{MSE})$ to compare
prediction stability; Fig.~4 varies noise frequency and the number of series.
In selected coefficient-recovery experiments, the model is also trained with
MAE loss to check sensitivity to the loss function. Figure~5 plots normalized
test loss against candidate input-window lengths; inspection of stable
coefficients at successive lags provides the complementary estimate of
effective memory depth.

For explanation recovery, the authors first average the sample-specific
$(\alpha^{\mathrm{DCI}}_t)_{k,j,\ell}$ over held-out windows within each run.
Let $\bar\alpha^{\mathrm{DCI},(r)}_{k,j,\ell}$ denote that reported coefficient
for run $r$. Across $R=5$ runs with different random seeds, paper
Eqs.~11--12 calculate

\begin{equation}
\begin{aligned}
    \mu_{k,j,\ell}
    &= \frac{1}{R}\sum_{r=1}^{R}
       \bar\alpha^{\mathrm{DCI},(r)}_{k,j,\ell}, \\
    \sigma_{k,j,\ell}
    &= \sqrt{\frac{1}{R-1}\sum_{r=1}^{R}
       \left(
       \bar\alpha^{\mathrm{DCI},(r)}_{k,j,\ell}
       -\mu_{k,j,\ell}
       \right)^2}.
\end{aligned}
\end{equation}

Their visualization retains a coefficient when

\begin{equation}
    I^{\mathrm{stable}}_{k,j,\ell}
    =\mathbb{I}\!\left[
        |\mu_{k,j,\ell}|>1.95\,\sigma_{k,j,\ell}
      \right]=1.
\end{equation}

This descriptive reproducibility filter retains coefficients whose across-run
mean dominates their variability. Ground-truth accuracy is assessed separately
by comparing the retained values or support with the generator. The official
code additionally applies
$|\mu_{k,j,\ell}|\geq 0.01$ when printing ``significant'' coefficients and
labels retained entries as present or absent in the nonzero ground-truth
support.\footnote{Official DCIts code:
\url{https://github.com/hc-xai/dcits}; see \texttt{src/utils.py}, especially
\texttt{collect\_multiple\_runs},
\texttt{calculate\_multiple\_run\_statistics}, and
\texttt{print\_significant\_alpha}.}
One implementation detail matters for exact replication: paper Eq.~12 uses the
sample standard deviation (denominator $R-1$), whereas the official code's
multiple-run statistics routine calls NumPy's default \texttt{std}, which uses
denominator $R$. This slightly changes the filter threshold for a fixed set
of runs.
The paper then compares mean learned coefficients with the exact generating
values in selected slices: Figs.~6--10 show lag-resolved patterns and
source-to-target summaries; the VAR(2) example prints each true lag-matrix
entry beside its learned mean $\pm$ standard deviation. In the cubic
experiment, Fig.~13 checks known negative linear and positive cubic
coefficients by polynomial branch. The lag-aggregated score
$\beta^{\mathrm{DCI}}_{t,k,j}$ is obtained from paper Eqs.~9--10 as

\begin{equation}
    \widetilde\beta^{\mathrm{DCI}}_{t,k,j}
    =\sum_{\ell=1}^{L}
      \left|(\alpha^{\mathrm{DCI}}_t)_{k,j,\ell}\right|,
    \qquad
    \beta^{\mathrm{DCI}}_{t,k,j}
    =\frac{\widetilde\beta^{\mathrm{DCI}}_{t,k,j}}
    {\sum_{j'=1}^{N}\widetilde\beta^{\mathrm{DCI}}_{t,k,j'}}.
\end{equation}

It tests which source series matters locally, but deliberately discards lag
and sign. For reported source-to-target summaries, the implementation applies
the same equations to $\alpha^{\mathrm{DCI}}$ after averaging it over held-out
windows, at which point the index $t$ is suppressed. The published explanation
assessment combines these numerical entrywise comparisons and plots. A single
coefficient-tensor MAE/RMSE, support F1, sign-accuracy, or local-contribution
error would extend the reported benchmark metrics.

